\subsection{Applying the Factor Theorem}
\begin{Example}
Write a polynomial with roots $5$, $-3$, $0$ and $4$.
\begin{lpic}{}{6}
\end{lpic}
\end{Example}
\begin{itemize}
%\item We can add \makeblank[1.5in] to the factors if we want to without changing the roots, but that gives a different polynomial.
\item We can only add other factors if we're sure they're never \makeblank[1.5in]
\item We can be sure that a \makeblank[1.5in] is never \makeblank[1.5in]
\end{itemize}
\begin{Example}
Write a polynomial with roots $-3$, $4$, and $\sqrt{5}$ and leading coefficient 7.
\begin{lpic}{}{5}
\end{lpic}
\end{Example}
\begin{itemize}
\item We can leave these polynomials in \makeblank[1.5in] form unless we need another form.
\end{itemize}
\begin{Example}
Write a polynomial with roots $-4$, $\sqrt{3}$, and $-\sqrt{3}$ and leading coefficint 1 in standard form.
\begin{lpic}{}{9}
\regtextleft{1}{-9}{The polynomial is};
\end{lpic}
\begin{itemize}
\item We had to multiply this out because \makeblank[3in]
\end{itemize}
\end{Example}